\section{第一天}
\begin{enumerate}
  \item 什么是\bs{依法治国}。
        \begin{enumerate}
          \item 依法治国是指广大人民群众在\zx{党的领导}下，\zx{依照宪法和法律}的规定，
          \item \zx{通过各种途径}和形式\zx{管理国家}事务，管理\zx{经济文化}事业，管理\zx{社会}事务，
          \item \zx{保障}国家各项工作都\zx{依法进行}，
          \item \zx{逐步实现社会主义民主的制度化、法律化}，
          \item 使这种\zx{制度和法律不因}领导人的改变而改变，不因领导人看法和注意力的改变而改变。
          \item \zd{依法治国}作为\zx{党领导人民治理国家}的\zx{基本方略}，
          \item 其\zx{根本目的}在于\zx{保障人民}充分行使当家作主的\zx{权利}，维护人民当家作主的\zx{地位}。
        \end{enumerate}
  \item 依法治国的价值意义。
        \begin{enumerate}
          \item \blh 是\zx{坚持和发展}中国特色社会主义的\zx{本质要求和重要保障}。
          \item \blh 是\zx{实现}\zx{国家治理体系}和治理\zx{能力现代化}的\zx{必然}要求。
          \item \blh \zx{事关}我们\zx{党执政兴国}、事关\zx{人民幸福安康}、事关\zx{党和国家长治久安}。
          \item \blh 也是\zx{全面推进中国式现代化建设}、\zx{实现}中华民族伟大\zx{复兴}的中国\zx{梦}\scx{的必然要求}。
          \item \blh 是\zx{全面深化改革、完善和发展}中国特色社会主义\zx{制度}、\zx{提高党}的\zx{执政能力和执政水平}\scx{的应有之意}。
        \end{enumerate}
  \item 【老马识途】全面依法治国的价值，可以理解为是\bs{万能价值}，
        但凡与法治相关的主题，都可以移花接木，
        照葫芦画瓢将对依法治国的价值意义的概括移植到\zd{对其他相关主题的价值阐述上}。典型如：
        \begin{enumerate}
          \item \zd{深化司法体制改革，推进实现司法公正}，
                是坚持和发展中国特色社会主义的本质要求和重要保障，
                是实现国家治理体系和治理能力现代化的必然要求，
                事关我党执政兴国，
                事关人民幸福安康，
                事关党和国家长治久安。
                全面推进中国式现代化建设、实现中华民族伟大复兴的中国梦，
                全面深化改革、完善和发展中国特色社会主义制度，
                提高党的执政能力和执政水平，
                \zd{必须全面推进司法体制改革，实现司法公正}。
          \item \zd{建设法治政府，推进政府职能转变}，
                是坚持和发展中国特色社会主义的本质要求和重要保障，
                是实现国家治理体系和治理能力现代化的必然要求，
                事关我党执政兴国，
                事关人民幸福安康，
                事关党和国家长治久安。
                全面推进中国式现代化建设、实现中华民族伟大复兴的中国梦，
                全面深化改革、完善和发展中国特色社会主义制度，
                提高党的执政能力和执政水平，
                \zd{必须全面推进政府职能转变，建设法治政府}。
        \end{enumerate}
\end{enumerate}


